
% \section*{Theoretical Justification for Sensitivity-Constrained Training} \label{prove}

% We provide a mathematical basis for why Sensitivity-Constrained Neural Operators (SC-NOs) outperform standard training by incorporating sensitivity loss, improving generalization, long-term stability, and inversion accuracy.

% We define the operator $\mathcal{G}_\theta: \mathcal{U} \to \mathcal{V}$, which maps an input function $a(x) = (u_0(x), p(x))$ to $\hat{u}(x,t) = \mathcal{G}_\theta(a(x))$, approximating the true solution $u(x,t)$.

% \subsubsection*{Assumptions and Notation}

% We work in bounded domains: the parameter space $\mathcal{P}_{\text{param}} \subset \mathbb{R}^m$ and the spatio-temporal domain $\mathcal{X} \times [0,T]$. We assume $\hat{u}$ is differentiable with respect to $p$. We use the Euclidean norm $|\cdot|$ on $\mathbb{R}^m$ and the corresponding induced operator norms.

% We assume the true solution $u$ and its sensitivity $\frac{\partial u}{\partial p}$ are bounded by $M$ and $M_p$ respectively, for all $(x,t,p)$.
% \begin{equation}
% |u(x,t;p)| \leq M, \quad \left\| \frac{\partial u}{\partial p}(x,t;p) \right\| \leq M_p.
% \end{equation}

% Our core argument relies on standard (though strong) assumptions from learning theory: for an overparameterized, well-trained model $\theta_s^*$, a small empirical loss implies a small uniform error bound.

% \textbf{Assumption 1 (Uniform Convergence of Gradients):}
% We assume that a model $\theta_s^*$ trained to $L_s(\theta_s^*)$ on $n$ samples satisfies:
% \begin{equation}
% \sup_{x,t,p} \left\| \frac{\partial \hat{u}}{\partial p}(x,t;p;\theta_s^*) - \frac{\partial u}{\partial p}(x,t;p) \right\| \le \varepsilon_s
% \end{equation}
% where $\varepsilon_s \to 0$ as $L_s \to 0$ and $n \to \infty$.

% \textbf{Assumption 2 (Uniform Convergence of Solutions):}
% We assume a similar uniform bound $\varepsilon_u$ for the solution itself:
% \begin{equation}
% \sup_{x,t,p} |\hat{u}(x,t;p;\theta_s^*) - u(x,t;p)| \le \varepsilon_u
% \end{equation}
% where $\varepsilon_u \to 0$ as $L_u \to 0$ and $n \to \infty$.

% \subsubsection*{Cost Functionals}

% Training minimizes the following losses, where $\frac{\partial u}{\partial p}$ is computed on-the-fly via a differentiable solver or adjoint model:
% \begin{equation}
% L_u(\theta) = \frac{1}{n} \sum_{i=1}^n [\hat{u}(x_i, t_i; p_i;\theta) - u(x_i, t_i; p_i)]^2,
% \end{equation}
% \begin{equation}
% L_s(\theta) = \frac{1}{n} \sum_{i=1}^n \left\| \frac{\partial \hat{u}}{\partial p}(x_i, t_i; p_i;\theta) - \frac{\partial u}{\partial p}(x_i, t_i; p_i) \right\|^2,
% \end{equation}
% \begin{equation}
% L(\theta) = L_u(\theta) + \lambda L_s(\theta), \quad \lambda > 0.
% \end{equation}


% \subsubsection*{Improved Generalization}

% Training with $L(\theta)$ acts as a powerful regularizer by explicitly constraining the model's Lipschitz constant with respect to parameters. We define this constant as:
% \begin{equation}
% L_{\text{lip}}(\theta) = \sup_{x,t,p} \left\| \frac{\partial \hat{u}}{\partial p}(x,t;p;\theta) \right\|.
% \end{equation}
% For the sensitivity-constrained model $\theta_s^*$, we can rigorously bound its Lipschitz constant using Assumption 1 and the triangle inequality:
% \begin{align}
% L_{\text{lip}}(\theta_s^*) &= \sup \left\| \frac{\partial u}{\partial p} + \left( \frac{\partial \hat{u}}{\partial p} - \frac{\partial u}{\partial p} \right) \right\| \notag \\
% &\le \sup \left\| \frac{\partial u}{\partial p} \right\| + \sup \left\| \frac{\partial \hat{u}}{\partial p} - \frac{\partial u}{\partial p} \right\| \notag \\
% &\le M_p + \varepsilon_s.
% \end{align}
% By driving $\varepsilon_s \to 0$, we force the model to inherit the (bounded) Lipschitz constant of the true physical system.

% This has a direct impact on generalization. For fixed $(x,t)$, we define the hypothesis class $\mathcal{F}_{\text{SC}} = \{p \mapsto \hat{u}(x,t;p;\theta): L_{\text{lip}}(\theta) \le M_p+\varepsilon_s\}$. On a bounded parameter domain, standard results imply that the empirical Rademacher complexity $\hat{R}_n(\mathcal{F}_{\text{SC}})$ of this class scales linearly with its Lipschitz constant, $\hat{R}_n(\mathcal{F}_{\text{SC}}) \lesssim (M_p + \varepsilon_s)$. This complexity is strictly lower than that of an unconstrained class $\mathcal{F}_u$ that admits arbitrarily large Lipschitz constants. The standard generalization bound,
% \begin{equation}
% \mathbb{E}[(\hat{u} - u)^2] \leq L_u(\theta) + 2 \hat{R}_n(\mathcal{F}) + \text{ComplexityPenalty}(\delta, n),
% \end{equation}
% is therefore tighter for $\mathcal{F}_{\text{SC}}$, implying better performance on unseen parameters.


% \subsubsection*{Long-Term Prediction Stability}

% This argument is most clearly made by setting the parameter $p$ to be the initial condition $u_0$. The sensitivity $\frac{\partial u(t)}{\partial u_0}$ is the \textbf{fundamental solution operator}, and its norm governs the growth of perturbations (i.e., the Lyapunov exponents).

% Assume the true physical system's growth is bounded:
% \begin{equation}
% \left\| \frac{\partial u(t)}{\partial u_0} \right\| \le C e^{\alpha t}
% \end{equation}
% for some physical exponent $\alpha$ (which could be positive for a chaotic system, or negative for a stable one).

% A standard model $\theta_u^*$, trained only on $L_u$, may learn an incorrect operator $\frac{\partial \hat{u}}{\partial u_0}$ with a much larger exponent $\hat{\alpha} \gg \alpha$, leading to spurious amplification of small errors.

% The SC-NO, by minimizing $L_s$ (with $p=u_0$), uses Assumption 1 to force its learned operator to be close to the true one:
% \begin{align}
% \left\| \frac{\partial \hat{u}(t)}{\partial u_0} \right\| &\le \left\| \frac{\partial u(t)}{\partial u_0} \right\| + \sup \left\| \frac{\partial \hat{u}}{\partial u_0} - \frac{\partial u}{\partial u_0} \right\| \notag \\
% &\le C e^{\alpha t} + \varepsilon_s.
% \end{align}
% This shows that the SC-NO reproduces the physical growth rates of perturbations up to $O(\varepsilon_s)$ and does not introduce spurious instabilities beyond those present in the governing system.


% \subsubsection*{Inversion Accuracy}

% In parameter inversion, we seek to find $p^*$ by minimizing a loss $\ell$ between our model $\hat{u}$ and an observation $u_{\text{obs}}$. Let $J(p) = \ell(u(p), u_{\text{obs}})$ be the "true" (but intractable) loss landscape, and $J_\theta(p) = \ell(\hat{u}_\theta(p), u_{\text{obs}})$ be the surrogate landscape our optimizer actually sees.

% Gradient descent updates use the gradient of $J_\theta$:
% \begin{equation}
% p_{k+1} = p_k - \eta \nabla_p J_\theta(p_k), \quad \text{where} \quad \nabla_p J_\theta = \frac{\partial \ell}{\partial \hat{u}} \cdot \frac{\partial \hat{u}}{\partial p}.
% \end{equation}
% The success of the inversion depends on the \textbf{gradient error} $\|\nabla J_\theta(p) - \nabla J(p)\|$. If $\ell$ is smooth, this error is bounded by our assumed uniform errors:
% \begin{equation}
% \|\nabla J_\theta(p) - \nabla J(p)\| \le C_1 \varepsilon_u + C_2 \varepsilon_s
% \end{equation}
% for some constants $C_1, C_2$.

% A standard model $\theta_u^*$ is only trained to minimize $\varepsilon_u$, leaving a large, uncontrolled gradient error $C_2 \varepsilon_s$. The optimizer is fed an inaccurate gradient. An SC-NO is explicitly trained to drive both $\varepsilon_u \to 0$ and $\varepsilon_s \to 0$. We further assume that, in a neighborhood of $p^*$, the true objective $J(p)$ is $\mu$-strongly convex. Under this condition, standard perturbation results for gradient descent imply that optimization on $J_\theta$ converges to an $O(\varepsilon_u + \varepsilon_s)$ neighborhood of the true optimum $p^*$. By minimizing $\varepsilon_s$, SC-NOs guarantee a high-fidelity gradient, ensuring a much more accurate and stable convergence.


















% \section*{Theoretical Justification for Sensitivity-Constrained Training} \label{prove}

% We provide a mathematical basis for why Sensitivity-Constrained Neural Operators (SC-NOs) outperform standard training by incorporating sensitivity loss, improving generalization, long-term stability, and inversion accuracy.

% We define the operator \(\mathcal{G}_\theta: \mathcal{U} \to \mathcal{V}\), which maps \( a(x) = (u_0(x), p(x)) \) to \(\hat{u}(x,t) = \mathcal{G}_\theta(a(x))\), approximating the true solution \( u(x,t) \). Both are continuously differentiable with respect to \( p \in \mathbb{R}^m \) over a bounded domain \(\mathcal{X} \times [0,T]\). The dataset consists of i.i.d. samples 

% \begin{equation}
% D = \{ (x_i, t_i, p_i, u(x_i, t_i; p_i), \frac{\partial u}{\partial p}(x_i, t_i; p_i)) \}_{i=1}^n
% \end{equation}

% drawn from a distribution \(\mathcal{P}\) over \(\mathcal{X} \times [0,T] \times \mathcal{P}_{\text{param}}\), where \(\mathcal{P}_{\text{param}}\) is a bounded subset of \(\mathbb{R}^m\) with \( \|p\| \leq P_{\text{max}} \). 

% The function outputs are bounded, with constants \( M, M_p > 0 \) satisfying

% \begin{equation}
% |\hat{u}(x,t;p;\theta)|, |u(x,t;p)| \leq M, \quad \left\| \frac{\partial \hat{u}}{\partial p}(x,t;p;\theta) \right\|, \left\| \frac{\partial u}{\partial p}(x,t;p) \right\| \leq M_p.
% \end{equation}

% The optimization landscape is assumed to be well-posed, with practical optima \(\theta_u^*\) and \(\theta_s^*\) achieving near-global minima in \(\Theta\). Training minimizes the following losses:

% \begin{equation}
% L_u(\theta) = \frac{1}{n} \sum_{i=1}^n [\hat{u}(x_i, t_i; p_i;\theta) - u(x_i, t_i; p_i)]^2,
% \end{equation}

% \begin{equation}
% L_s(\theta) = \frac{1}{n} \sum_{i=1}^n \left\| \frac{\partial \hat{u}}{\partial p}(x_i, t_i; p_i;\theta) - \frac{\partial u}{\partial p}(x_i, t_i; p_i) \right\|^2,
% \end{equation}

% \begin{equation}
% L(\theta) = L_u(\theta) + \lambda L_s(\theta), \quad \lambda > 0.
% \end{equation}

% \subsubsection*{Improved Generalization}

% For sufficiently large \(\lambda\), training with \(L(\theta)\) reduces complexity and improves generalization compared to \(L_u(\theta)\) alone. The Lipschitz constant of the model with respect to \( p \) is defined as:

% \begin{equation}
% L_{\text{lip}}(\theta) = \sup_{x,t,p} \left\| \frac{\partial \hat{u}}{\partial p}(x,t;p;\theta) \right\|.
% \end{equation}

% For any \( p, q \in \mathcal{P}_{\text{param}} \),

% \begin{equation}
% |\hat{u}(x,t;p;\theta) - \hat{u}(x,t;q;\theta)| \leq L_{\text{lip}}(\theta) \|p - q\|.
% \end{equation}

% With sensitivity training,

% \begin{equation}
% L_{\text{lip}}(\theta_s^*) \leq M_p + \sqrt{n L_s(\theta_s^*)},
% \end{equation}

% which ensures \( L_{\text{lip}}(\theta_s^*) < L_{\text{lip}}(\theta_u^*) \) as \( \lambda \to \infty \). This reduction in complexity leads to a lower Rademacher complexity,

% \begin{equation}
% \hat{R}_n(\mathcal{F}) \leq 2 L_{\text{lip}}(\theta) D \sqrt{\frac{2 \ln(2)}{n}}.
% \end{equation}

% Thus, the generalization bound

% \begin{equation}
% \mathbb{E}[(\hat{u} - u)^2] \leq L_u(\theta) + 2 \hat{R}_n(\mathcal{F}) + 3 \sqrt{\frac{\ln(2/\delta)}{2n}}
% \end{equation}

% is tighter for sensitivity-trained models, improving overall generalization.

% \subsubsection*{Long-Term Prediction Stability}

% Training with \( L(\theta) \) reduces error accumulation over time in dynamical systems. Defining the prediction error,

% \begin{equation}
% e(x,t;p;\theta) = \hat{u}(x,t;p;\theta) - u(x,t;p),
% \end{equation}

% the error growth is bounded by

% \begin{equation}
% |e(x,t;p;\theta)| \leq \int_0^t \left\| \frac{\partial \hat{u}}{\partial p}(x,s;p;\theta) - \frac{\partial u}{\partial p}(x,s;p) \right\| \|p\| \, ds.
% \end{equation}

% With probability \( 1 - \delta \),

% \begin{equation}
% \left\| \frac{\partial \hat{u}}{\partial p}(x,s;p;\theta_s^*) - \frac{\partial u}{\partial p}(x,s;p) \right\| \leq \sqrt{\frac{n L_s(\theta_s^*)}{\delta}},
% \end{equation}

% leading to

% \begin{equation}
% |e(x,t;p;\theta_s^*)| \leq t \sqrt{\frac{n L_s(\theta_s^*)}{\delta}} P_{\text{max}}.
% \end{equation}

% Since \( L_s(\theta_s^*) \to 0 \) as \( \lambda \to \infty \), sensitivity training mitigates long-term prediction errors.

% \subsubsection*{Inversion Accuracy}

% In parameter inversion, the goal is to solve

% \begin{equation}
% p^* = \arg\min_p \ell(\hat{u}(x,t;p;\theta), u_{\text{obs}}).
% \end{equation}

% Gradient updates follow

% \begin{equation}
% p_{k+1} = p_k - \eta \nabla_p \ell(p_k), \quad \nabla_p \ell = \frac{\partial \ell}{\partial \hat{u}} \cdot \frac{\partial \hat{u}}{\partial p}.
% \end{equation}

% For sensitivity-trained models, with probability \( 1 - \delta \),

% \begin{equation}
% \left\| \frac{\partial \hat{u}}{\partial p}(x,t;p;\theta_s^*) - \frac{\partial u}{\partial p}(x,t;p) \right\| \leq \sqrt{\frac{n L_s(\theta_s^*)}{\delta}},
% \end{equation}

% which ensures more accurate gradients and improved parameter recovery.\\
% We have demonstrated that training with \( L(\theta) = L_u(\theta) + \lambda L_s(\theta) \) provides key advantages. By constraining the Lipschitz constant, SC-NOs reduce complexity and achieve tighter generalization bounds. Minimizing sensitivity error stabilizes long-term predictions, and enforcing accurate gradients enhances parameter inversion. These results validate the effectiveness of SC-NOs in learning physically consistent representations, reinforcing their advantages over standard Neural Operators.
